\subsection{محدودیت‌ها و تحلیل شکست}
تحلیل داده‌های پایگاه اعتبارسنجی و گزارش‌های ثبت خطا در جدول \ref{tab:benchmark_replace} نکات زیر را به وضوح تبیین می‌نماید:
\begin{itemize}
    \item \textbf{علت وقوع دقت صفر درصد (\lr{Accuracy=0.00\%}) و وضعیت \lr{FAIL}:} در هر سه مدل ارزیابی‌شده (\lr{DistilBERT}، \lr{RoBERTa} و \lr{GPT-2})، دقت ثبت‌شده دقیقاً \lr{0.00\%} و وضعیت نهایی به صورت \lr{FAIL} گزارش گردیده است. بررسی مستقیم فایل‌های گزارش و لاگ اجرای آزمایش‌ها (\lr{summary.csv}) نشان می‌دهد که علت شکست، نقص معنایی یا ناتوانی محاسباتی ماژول نبوده، بلکه ناشی از خطای تزریق معماری در زمان اجرای بنچمارک تحت عنوان زیر است:
    \begin{center}
    \lr{\texttt{RuntimeError: Could not inject attention 'replace': No valid attention class found in module 'replace'}}
    \end{center}
    این خطای تزریق گراف محاسباتی زمانی رخ داده است که تابع جستجوی ثبت‌کننده ماژول‌ها (\lr{ATTENTION\_REGISTRY}) نتوانسته است کلاس هدف را مستقیماً از نام ماژول استخراج کند. در نتیجه ارزیابی مدل با لایه توجه خالی یا نامعتبر اجرا شده و ارزیاب خط لوله به منظور پایبندی به قوانین تضمین کیفیت، امتیاز دقت را صفر لحاظ کرده است.
    \item \textbf{حفظ شاخص‌های آماری خط پایه:} همان‌طور که در جدول مشاهده می‌شود، مقادیر شاخص F1 به ترتیب برابر با \lr{0.9077}، \lr{0.7586} و \lr{0.6354} با حافظه و توان عملیاتی خط پایه ثبت شده‌اند که تاییدی بر خوانش آماری پیش از شکست تزریق لایه است.
\end{itemize}
